\section{The One Who Does Not Receive}

\subsection{Not a being that happens to exist}

Finite things do not exhaust existence. A human being exists humanly; an oak exists according to the bounded nature of an oak; an angel, in Thomistic theology, would exist immaterially yet still according to a received created nature. In each case the answer to \emph{what is it?} does not by itself explain \emph{that it is}. Its act of being is possessed under a limit.

Aquinas gives this distinction its most concentrated form. Where essence and existence are distinct, existence is caused; a thing cannot give itself the act by which it is available to act at all. The first efficient cause therefore cannot be a being whose existence is received. In God alone, essence is not one principle and existence another. God does not merely \emph{have} being. He is his act of being (ST I, q.~3, a.~4; \emph{De ente et essentia}, ch.~4).

The traditional formula is \latin{ipsum esse subsistens}: subsistent being itself. It must be heard carefully. God is not ``being in general,'' a thin abstraction formed by noticing what creatures share. Nor is he a universal substance of which creatures are pieces, a cosmic field, or the sum total of everything. Common being is a concept predicable of many. Subsistent being is the one unreceived actuality from whom the many receive their diverse acts of existence (ST I, q.~4, a.~2; q.~44, a.~1).

This is why the regress does not discover a hidden object behind visible objects. It discovers that the source of every object is not an object in the creaturely sense. God and the world do not add up to two specimens within a larger inventory called reality. There is no larger field containing them both. The Creator--creature distinction is more radical: one is being without reception; every other reality has being by participation.

\subsection{Without seams}

Anything composed depends upon principles that account for its unity. Matter receives form. A whole requires its parts to be together. An accident modifies a subject. A member of a genus is specified and limited by what distinguishes it from other members. If the first source were composite in any of these ways, the unity and actuality of its components would still call for explanation.

Pure act is therefore simple---not easy, featureless, or impoverished, but without metaphysical seams. In God there is no composition of matter and form, essence and existence, substance and accident, genus and difference, or unrealized capacity and acquired act (ST I, q.~3). Wisdom, goodness, power, life, and love are not detachable properties fastened to a neutral divine core. They name from different creaturely angles the one infinite actuality that is God.

Divine simplicity does not shrink God into blankness. It means that every pure perfection found divided and limited in creatures pre-exists in the first source without division or limit. Because God is subsistent being, no perfection of being is absent from him; what is dispersed among creatures is one in its uncreated cause (ST I, q.~4, a.~2). The One is not less than the many. The many are fragmentary likenesses of the fullness they cannot contain.

\subsection{Why there can be only one}

Two numerically distinct divine essences or independent first sources would have to differ. One would possess something the other lacked, or each would instantiate a common nature under some differentiating condition. Either way, at least one would be limited by received or composite actuality. But subsistent being is not received into a limiting nature, and pure act lacks no perfection of being. Aquinas therefore argues from simplicity and infinite perfection that there cannot be several gods. This establishes unity of divine essence; it does not exclude the revealed distinction of Persons by subsistent relations of origin (ST I, q.~11, a.~3; q.~29, a.~4).

The unity reached here is not the number one set beside the number two. It is the indivisibility of the unconditioned source. God is not one because competitors happen not to exist. He is one because the actuality that is wholly unreceived and unrestricted cannot be multiplied by receivers, matter, differences, or limits (ST I, q.~11, a.~4).

Christian tradition hears a profound consonance between this metaphysical judgment and the name revealed to Moses: ``I AM'' (Exod 3:14). The scriptural name is not a scholastic formula smuggled into ancient Hebrew, and the proof does not rest on a translation of one verse. Yet Aquinas and the wider tradition rightly receive the name as luminous: creatures can say ``I am this'' or ``I am now''; the One without origin says simply that he is (ST I, q.~13, a.~11; CCC 203--213).

Try to let the negations have their weight. No cause behind him. No material beneath him. No possibility waiting within him. No future in which he becomes more fully himself. No larger reality against which he is measured. No second source dividing the field. The mind has traced every dependent ``because'' to the One in whom essence and existence are not two.

At that point one does not yet comprehend God. One has learned why comprehension must fail. There is nowhere outside the source of being from which to take his measure.
